%!TEX root = ../../main.tex
\section{교전 집합을 정하는 상수}
\label{sec:eng-constants}

표 \ref{tab:constants}는 교전 집합과 규모 계급과 결과를 정하는 상수를 값·출처와 함께 모은 것이다. 출처 칸은 네 가지이다. 자료 추정은 코퍼스에서 얻은 값이며 추정값과 구간을 갖는다. 게임 규칙은 게임이 쓰는 값에 맞춘 가정이며 패치 노트로 검사한다. 관례는 발표된 정의나 이전 연구를 따른 값이다. 구현 선택은 구현하면서 정한 값이며, 측정된 영향을 적거나 재지 않았음을 적는다. 값은 검출기에 넣은 값이고, 시간 경계와 거리 경계는 넣기 전에 반올림하였다. 이 표는 요약이며, 전체 열아홉 행과 행마다의 근거는 부록 \ref{app:constants}의 표 \ref{tab:app-constants}에 있다.

%% table 객체: tab:constants — 요약판. 출처 Methods §1 상수표 (sec_definition.tex tab:constants L451–L641); 전체는 부록 B.
%% 2026-09-22: 표주 제거(본문으로), 출처 분류 이름을 본문과 통일(자료 추정 / 게임 규칙 / 관례 / 구현 선택), 기호 τ와 코드 세부 제거.
\begin{table}[!t]
  \centering
  \caption{교전 집합을 정하는 주요 상수}
  \label{tab:constants}
  \small
  \begin{tabularx}{\linewidth}{@{}l>{\raggedright\arraybackslash}p{2.6cm}>{\raggedright\arraybackslash}p{1.6cm}X@{}}
    \toprule
    상수 & 값 & 출처 & 근거와 측정된 영향 \\
    \midrule
    시간 경계 $G$ & \constG{} s (추정 \constGest{}) & 자료 추정 & 킬 간격 분포의 두 봉우리 사이; 95\% 구간 \constGciLo{}--\constGciHi{} s; \ariPlateauLo{}--\ariPlateauHi{} s에서 묶음이 유지됨 \\
    거리 경계 $D$ & \constD{} u (추정 \constDest{}) & 자료 추정 & 챔피언 공유율이 0.5로 떨어지는 거리; 95\% 구간 \constDciLo{}--\constDciHi{} u (표본 변동만) \\
    존재 반경 $R$ & \constR{} u & 게임 규칙 & 사망 경험치 공유 반경; 패치 노트에 변경 없음 \\
    여유 $B$ & \constB{} s & 게임 규칙 & 어시스트 인정 시간; 예측 시점을 정함 \\
    팀당 최소 $M$ & \constM{} (양 팀, $R$ 안, 생존) & 관례 & Ke 등을 따름; 3이면 \numEngMthree{} 교전 \\
    위치 격자 & 5 s; 프레임·킬 궤적 보간 & 구현 선택 & 파일럿 중앙 위치 오차 1{,}149 u; 존재 조건에 대한 영향 미측정 \\
    재교전 병합 & 15 s / 2{,}000 u & 구현 선택 & 반경을 1{,}000--3{,}000 u로 바꾸면 교전 수가 1\% 안에서 움직이고 계급은 옮기지 않음 \\
    지속 시간 상한 & 예측 시점 이후 60 s & 구현 선택 & 45 s와 90 s에서 교전 수가 1\% 남짓 줄고 계급은 옮기지 않음 \\
    시작·끝 제외 & 첫 2분; 마지막 프레임 35 s 전 & 구현 선택 & 미측정 \\
    관측 창 & 예측 시점까지 30 s, 5 s 구간 & 관례 (이전 정의) & 창 길이는 이전 연구의 라벨 아래에서만 훑음 \\
    상호작용 반경 & 3{,}000 u; 상점 사건 포함 & 구현 선택 & 참여 인원, 곧 규모 계급을 정함; 이전 정의에서 이어받았고 값의 근거는 기록 없음; 2{,}000--4{,}264 u에서 교전 수는 그대로이고 한타가 \oatItwoT{}--\oatIfourdT{}개로 움직임 \\
    규모 절단 & pick $\le 1$ / 소규모 2--3 / 한타 $\ge 4$ & 관례 & pick과 한타의 AUC 차이 부호가 절단에 따라 바뀜 \\
    \bottomrule
  \end{tabularx}
\end{table}


표에 없는 설정이 둘 있다. 검출기가 읽는 1{,}800 단위 근접 반경은 참가자별 진단에만 쓰고 교전 집합과 무관하다. 설정에 남아 있는 이전 검출기의 스위치들도 지금의 킬 기반 검출기가 읽지 않는다. 결과를 재는 끝 시각의 상수는 제 \ref{sec:eng-endpoint} 절에서 따로 정한다.
